Put \(m_i=\{y_i(1)+y_i(0)\}/2\) and \(\delta_i=y_i(1)-y_i(0)\).  The potential outcomes are fixed by Assumption \(\text{ass:fixed-potential-outcomes}\), so all expectations below are with respect to the assignment alone.  Assumption \(\text{ass:consistency-no-interference}\) and \(z_i\in\{-1,1\}\) give, identically for every supported assignment,
\[
Y_i^{\mathrm{obs}}=m_i+\frac{z_i\delta_i}{2}.
\]
Consequently,
\[
\widehat\tau_d
=\frac{2}{n}\sum_{i=1}^n z_i m_i+\frac{1}{n}\sum_{i=1}^n z_i^2\delta_i
=\tau_n+\frac{2}{n}z^\top m.
\]
Every \(z\in\mathcal Z_n\) is orthogonal to \(\mathbf 1_n\).  Since \(\mu=P_Hm\), this implies \(z^\top m=z^\top\mu\), and hence
\[
\widehat\tau_d-\tau_n=\frac{2}{n}z^\top\mu.
\]
The defining mean-zero condition in \(\mathcal D_n\) yields \(\mathbb E_d[z]=0\).  Thus \(\mathbb E_d[\widehat\tau_d]=\tau_n\).  It also implies \(A_d=\mathbb E_d[zz^\top]\), and therefore
\[
\mathbb E_d[(\widehat\tau_d-\tau_n)^2]
=\frac{4}{n^2}\mathbb E_d[\mu^\top zz^\top\mu]
=\frac{4}{n^2}\mu^\top A_d\mu.
\]
Finally, because each \(z_i\) is binary and has mean zero, \(\Pr_d(z_i=1)=\Pr_d(z_i=-1)=1/2\).  Thus the marginal exposure probabilities used by the unbiasedness calculation are positive; no superpopulation or outcome randomness has been introduced.